\section*{§1. Sets and Mappings}

\textbf{Definition 1.} A \textit{set} is a collection of objects, called the \textit{elements} of the set.

We write $x \in A$ if $x$ is an element of the set $A$, and $x \notin A$ if $x$ is not an element of the set $A$.

A set can be defined by listing its elements between braces. For example, the set $A$ consisting of the numbers 1, 2, and 3 is written as

\[
A = \{1, 2, 3\}.
\]

The order in which the elements are listed is irrelevant. Thus, $\{1, 2, 3\} = \{3, 1, 2\}$.

If a set has infinitely many elements, then we cannot list all of them. In this case, we can use an ellipsis to indicate that the list continues in an obvious way. For example, the set of all positive integers is written as

\[
\{1, 2, 3, \dots\}.
\]

Another way to define a set is to specify a property that the elements of the set must satisfy. For example, the set of all real numbers $x$ such that $x^2 < 2$ is written as

\[
\{x \in \mathbb{R} \mid x^2 < 2\}.
\]

The symbol $\mid$ is read as "such that."

\textbf{Definition 2.} If every element of a set $A$ is also an element of a set $B$, then we say that $A$ is a \textit{subset} of $B$, and we write $A \subset B$ or $B \supset A$.

If $A \subset B$ and $B \subset A$, then we say that $A$ and $B$ are \textit{equal}, and we write $A = B$.

If $A \subset B$ but $A \neq B$, then we say that $A$ is a \textit{proper subset} of $B$, and we write $A \subsetneq B$ or $B \supsetneq A$.

The \textit{empty set}, denoted by $\emptyset$, is the set that contains no elements. The empty set is a subset of every set.

\textbf{Definition 3.} The \textit{union} of two sets $A$ and $B$, denoted by $A \cup B$, is the set of all elements that are in $A$ or in $B$ (or in both).

The \textit{intersection} of two sets $A$ and $B$, denoted by $A \cap B$, is the set of all elements that are in both $A$ and $B$.

The \textit{difference} of two sets $A$ and $B$, denoted by $A \setminus B$, is the set of all elements that are in $A$ but not in $B$.

The \textit{complement} of a set $A$, denoted by $A^c$, is the set of all elements that are not in $A$. The complement of a set depends on the context. For example, if we are working with real numbers, then the complement of the set of all positive real numbers is the set of all non-positive real numbers.